\section{Stakeholders and needs}
In any project, identifying the stakeholders and understanding their needs is essential to ensuring the system delivers real value. For Quizum, stakeholders come from different backgrounds, including learners, educators, developers, and institutions. Each group has unique expectations and priorities that shape how the system should be designed, developed, and maintained.
\begin{itemize}
    \item \textbf{Students:} Students are the primary users of the system, and their main need is a platform that saves them time and enhances learning efficiency. They want video lectures summarized into concise, accurate, and easy-to-read text so that they can quickly revise or revisit key points without watching the entire lecture again. For them, reliability, accessibility across devices, and clarity of summaries are crucial to ensure that they truly benefit from the system in their daily study routines.
    \item \textbf{Educators:} they rely on the platform not only to support their teaching but also to improve student engagement. They need features that provide insights into how students interact with the summarized content, such as which parts of a lecture students review most often. This helps them identify knowledge gaps and refine their teaching strategies. Additionally, they value tools that ensure the accuracy of information so their teaching material is conveyed correctly and effectively.
    \item \textbf{Developers and technical teams:} they are critical stakeholders responsible for building and maintaining the system. They require a clear set of requirements, access to scalable technologies, and resources to implement robust solutions. Their needs revolve around ensuring system performance, security, and seamless integration with tools like cloud storage or learning management systems. They also need reliable feedback loops from users to continuously improve the platform.
    \item \textbf{Content creators:} such as lecturers who record video material, have needs that go beyond just student learning. They want assurance that their intellectual property is protected and not misused while being processed or summarized. At the same time, they need the system to make their lectures more impactful and accessible to students, increasing the overall effectiveness of their teaching without adding extra workload on their side.
\end{itemize}